% Introduction
\section{Introduction}

As our general knowledge about reality evolves, we tend to perceive erst hidden intricacy in things around us. Such epistemological evolution sometimes makes the seemingly simple again complicated, the supposedly solved yet again lacking an adequate solution; an improved comprehension oftentimes leads to more subtle questions about the phenomena in appraisal, which -- hopefully -- leads to more profound answers, in a dialectical process that results in an ever-increasing understanding about our reality.

Examples of such process and how it impacts our lives are rather abundant. Of particular interest is the nowadays widely known fact that `non-trivial' interrelations between even simple things might make their global behavior hard to understand -- let alone to predict --, sometimes even to the extent that said behavior seems to have an entirely new nature. Consider, for instance, the problem of consciousness, and how it leads us to believe, at least partially in an explanatory attempt, in things such as the existence of soul. The behavior that arises from what could ultimately be interacting matter being so wildly unexpected, it seems unreasonable to explain it without resourcing to things other than matter alone. 

The acknowledgment of the ubiquity of similar problems in the world that surrounds us has led to the development of formal sciences -- systems and complexity sciences -- that aim to study, among other things, in which ways an ensemble of entities can give rise to things as qualitatively different as matter seems to be from consciousness. Despite being quite new and yet under development, some hindrances dealt with by complexity science seem to be fundamental. One of them being the fact that complex systems tend to be hard to treat with some of the conventional mathematical tools used in closely-related problems, such as those used within dynamical systems theory. 

The inadequacy of some of these usual methods, together with the vigorous technological -- particularly, computational -- development we witnessed in the last decades, happened to elect computational simulations as one of the main tools for developing our knowledge about complex systems; and not unlike the way that some conventional methods started to fail us, some usual assumptions about our world also started to seem insufficient, to say the least. In response to that, we are also bound to witness an ever-increasing importance of a (within the bounds of reason) general theory of complex systems.

Most of the systems that are relevant to us are complex. Most of the hardest problems we face as society and species lie within the domain of complex systems theory. In order to venture forth into the future, it is paramount to be able to sturdily tackle complexity problems, which cannot be done without both the aid of computational science and a proper theoretical foundation for the understanding of such systems.

\subsection{Motivations and scope}

Distributed systems are ever more present in contemporary society. Added to that the increasing complexity of these systems, one can soundly defend the growing relevance of complex systems theory to enhance one's understanding of, and control over, these systems -- goals known as notably difficult. Among other concerns, those motivated the development of theories that intend to provide an adequate framework to design and understand these types of systems in all their intricacy. Of particular relevance to this work is that of \textit{holonic multiagent systems}. Here, nonetheless, we will have motivations other than those stated. Our interest lies primarily in developing tools to improve our understanding of complex systems in general, aiming to help in the search for the most general and formal theory possible for such systems. 

More specifically, we are particularly oriented towards understanding the idea of emergence and providing an adequate formalism to describe the emergence of properties and behaviors, despite the depth of the discussion this concept engenders being such that some few pages could hardly contain a framework capable of providing a thorough analysis. That notwithstanding, we believe this work might raise some useful insights, and perhaps outline a proficuous path of investigation. Since a core instrument in the study of complex systems is computer simulations, we believe that such a framework should be as tightly as possible related to computational models. Indeed, while developing the presented theoretical approach, it was a constant concern that it should be capable of both helping us to explain observed phenomena (either \textit{in vivo}, \textit{in vitro} or \textit{in silico}) and providing a convenient, structured scheme for devising models and, specially, their computational implementations. 

We are also particularly interested in its adequacy for the treatment of \textit{artificial life systems}, since those seem to provide a useful, virtually infinite, zoology for one to explore and develop insights about more general concepts concerning the theory of complex systems. In this sense, this work aims to more tightly integrate the interplay between \textit{understanding} and \textit{modeling}: we envision that the same theoretical tools we use to analyze properties of a system should be as immediately as possible applicable to simulate a system with said properties. 

\subsection{Evolution of scope and present contribution}

More properly about the work, we should say: our initial goal was to explore a kind of system model characterized by its nestedness and hierarchical structure, together with restrictions over the possible subsystems, in that they would be required to be `similar' to their corresponding supersystems, in a sense that is to be precised -- but not within this work. We would explore this restriction (together with other model `parameters', such as the very similarity metric) as a main value to be used both in analysis and modeling. It turned out that laying the theoretical foundations we thought convenient to describe and work with these systems in the way we intended was much more lengthy than we had anticipated. We thus had to postpone this initial goal to future works, and this one acquired an even more theoretical flavor than it was initially intended to have. This observation is pertinent since said original goal motivates many of the theoretical decisions we have made, and it will be occasionally adequate to mention it in order to provide context to a more sound justification of these choices.

As said, we are particularly interested in artificial life models. It is also true that we are undeniably biased towards the theory of dynamical systems. Although rather restrictive in scope -- and generally unfit to treat living systems --, it is an arguably well-established theory, with plenty of tools for both quantitative and qualitative analysis, that happens to have a significant proximity to the theory of complex systems, despite it yet not being entirely clear whether it stands as a supertheory, a subtheory, or none of them. That notwithstanding, the two theories describe time-evolving systems, so they are, at the very least, parts of an overarching theory-to-be. Approximating them as much as possible in order to transfer at least some of these instruments to the developing theory of complex systems would be as useful as this approximation permits.

\subsection{Guiding metaphor: spatial systems}

We are also heavily inspired by physical systems. The definitions we give are, in fact, abstractions based on the idea that any physical object -- in particular, system -- corresponds to a region in space (a supersubstantivalist commitment, while allowing interpenetration). The scale parameter has its most natural interpretation as a `spatial resolution' parameter, and, under this interpretation, varying it is understood as `zooming inwards/outwards' -- what explains the choice of some terms that will be presented later on. 

To what extent the developed framework transfers to broader contexts, and what adaptations, if any, are needed in order to do so, is not a question we will deal with for now, despite its obvious importance. The reason for that negligence is twofold: firstly, we believe that this theory is general enough so that it could be -- in the worst case scenario, with minor reworks -- easily transferable for other contexts, where systems do not exist in `space' (more precisely, in topological spaces); secondly, we believe that this framework is interesting enough even if not too immediately transferable to these other cases. For these reasons, we delegate the aforementioned investigation to future works.
